\section{Steiner Rooted Connected Orientation of Graphs and Hypergraphs}
\label{sec:rooted-steiner-connected-orientations}
In this section, we use  \Cref{thm:complete-splitting-off} to
prove Theorems \ref{thm:rooted-steiner-orientation-in-graphs} and \ref{thm:rooted-steiner-orientation-in-hypergraphs}. 
We recall the \textsc{Max Steiner Rooted-Connected Orientation} problem: the input here is a hypergraph $G=(V, E)$, a subset $T\subseteq V$ of terminals, and a root vertex $r\in T$. The goal is to find a maximum integer $k$ and an orientation $\overrightarrow{G}$ of $G$ such that $\overrightarrow{G}$ is Steiner rooted $k$-hyperarc-connected. 
%As warm-up to Theorems \ref{thm:rooted-steiner-orientation-in-graphs} and \ref{thm:rooted-steiner-orientation-in-hypergraphs}, we consider the 
Although \textsc{Max Steiner Rooted-Connected Orientation} is NP-hard, two extreme cases of the problem are polynomial-time solvable: the number of terminals $|T|=2$ and the number of terminals $|T|=|V|$. We discuss these polynomial-time solvable cases as warm-ups towards Theorems \ref{thm:rooted-steiner-orientation-in-graphs} and \ref{thm:rooted-steiner-orientation-in-hypergraphs}. The ideas underlying these cases will be useful in proving Theorems \ref{thm:rooted-steiner-orientation-in-graphs} and \ref{thm:rooted-steiner-orientation-in-hypergraphs}. 
We recall that non-terminals are denoted as Steiner vertices. 

\paragraph{The case of $|T|=2$.} \textsc{Max Steiner Rooted-Connected Orientation} problem where we have exactly $2$ terminals is polynomial-time solvable owing to Menger's theorem in hypergraphs. 
A path between $u$ and $v$ in a hypergraph is an alternative sequence of distinct vertices and hyperedges $v_1=u, e_1, v_2, e_2, ..., e_{\ell-1}, v_{\ell}=v$ such that $v_i, v_{i+1}\in e_i$ for every $i\in [\ell-1]$. 
%and $v_i\neq v_j$ for every distinct $i,j\in [\ell]$. 
We recall Menger's theorem: for a hypergraph $G=(V, E)$ and a pair of distinct vertices $r, v\in V$, there exist $k$ hyperedge-disjoint paths between $r$ and $v$ in $G$ if and only if $\lambda_G(r,v)\ge k$. We present a proof of this fundamental graph-theoretical result using \Cref{thm:complete-splitting-off}. Our proof also holds for undirected graphs. 
Parts of our proof will be useful in the proof of Theorems \ref{thm:rooted-steiner-orientation-in-graphs} and \ref{thm:rooted-steiner-orientation-in-hypergraphs}. 

Let $G=(V, E)$ be a hypergraph, $r,v\in V$ be distinct vertices, and $k$ be a positive integer. It is easy to see that if there exist $k$ hyperedge-disjoint paths between $r$ and $v$ in $G$, then $\lambda_G(r,v)\ge k$. We prove the converse. Suppose that $\lambda_G(r, v)\ge k$. We would like to show that there exist $k$ hyperedge-disjoint paths between $r$ and $v$ in $G$. 
Consider performing local connectivity preserving complete h-splitting-off operations at Steiner vertices in the current hypergraph and then deleting them. This will result in a $2$-vertex hypergraph $H=(\{r,v\}, E_H)$ such that $\lambda_H(r,v)\ge k$. Since $H$ contains only $2$ vertices, it immediately follows that all hyperedges of $H$ are in fact edges, and they can be oriented to obtain $k$ hyperarc-disjoint paths 
%from $r$ to $v$. 
from $v$ to $r$. 
Let $\overrightarrow{H}$ be such an orientation. 
We now extend this orientation $\overrightarrow{H}$ to an orientation $\overrightarrow{G}$ of the original hypergraph $G$ that still contains $k$ hyperarc-disjoint paths 
%from $r$ to $v$. 
from $v$ to $r$. 
We use Lemma \ref{lem:lifting-Steiner-orientation-through-splitting-off} below for this extension. 
Lemma \ref{lem:lifting-Steiner-orientation-through-splitting-off} below proves a more general statement for arbitrary number of terminals and for a single h-splitting-off operation. We recall that a complete h-splitting-off at a vertex from a hypergraph is a sequence of h-splitting-off operations. 
Lemma \ref{lem:lifting-Steiner-orientation-through-splitting-off} can be applied inductively starting from $\overrightarrow{H}$ to arrive at an orientation $\overrightarrow{G}$ of the hypergraph $G$ such that there exist $k$ hyperarc-disjoint paths 
%from $r$ to $v$ 
from $v$ to $r$ 
in $\overrightarrow{G}$. Such an orientation immediately gives the $k$ hyperedge-disjoint paths between $r$ and $v$ in $G$. This concludes our proof of Menger's theorem in hypergraphs and graphs (undirected). 
%We note that the proof of \Cref{thm:complete-splitting-off} does not rely on Menger's theorem, so this is a self-complete proof. 


\begin{lemma}\label{lem:lifting-Steiner-orientation-through-splitting-off}
    Let $G=(V, E_G)$ be a hypergraph, $T\subseteq V$ be a set of terminals, $r\in T$ be a specified root vertex, and $s\in V-T$. 
    Let $H = (V, E_H)$ be a hypergraph obtained by applying a h-splitting-off operation at the vertex $s$ from $G$. 
    Suppose there exists an orientation $\overrightarrow{H}$  
    %= (V, E_G, \head_G:E\rightarrow V)$ 
    of $H$ such that $\overrightarrow{H}$ is Steiner rooted $k$-hyperarc-connected. 
    Then, there exists an orientation 
    $\overrightarrow{G}$ 
    %= (V, E_H, \head_H:E_H\rightarrow V)$ 
    of $G$ such that $\overrightarrow{G}$ is Steiner rooted $k$-hyperarc-connected. 
\end{lemma}
\begin{proof}
 Let $\overrightarrow{H}= (V, E_H, \head_H:E_H\rightarrow V)$ be an orientation of $H$ that is Steiner rooted $k$-hyperarc-connected.  
  We will use the orientation $\overrightarrow{H}$ to define an orientation $\overrightarrow{G}=(V, E_G, \head_G: E_G\rightarrow V)$ for the hypergraph $G$. 
  We recall that $H$ is obtained by applying a h-splitting-off operation at $s$ from $G$. We have two cases: 
  %In particular, we define $\head_{H}:E_H\rightarrow V$ in two cases as follows:
\begin{enumerate}
    \item Suppose that $H$ is obtained 
    from $G$ by a h-merge almost disjoint hyperedges operation. Let $e, f \in \delta_{G}(s)$ be the hyperedges with $e\cap f = \{s\}$ that are h-merged during this operation. For notational convenience, we denote $g:= e\cup f$ as the hyperedge formed after the h-merge operation. Then, we define the orientation $\head_{G}:E_G\rightarrow V$ as follows: for each $e' \in E_G$,
$$\head_{G}(e') := \begin{cases}
    s& \text{ if } e' \in \{e, f\} \text{ and } \head_{H}(g) \not \in e'\\
    
    \head_{H}(g) &\text{ if } e' \in \{e, f\} \text{ and } \head_{H}(g) \in e'\\
    
    \head_{H}(e')& \text{ otherwise (i.e., $e'\in E_G-\{e, f\}$).}
\end{cases}$$
    %\knote{Replace ``otherwise'' above by ``$e'\in E_H-\{e, f\}$''}

    \item Suppose that $H$ is obtained from $G$ 
    %by a h-splitting-off at $s$ operation is a 
    by a h-trim hyperedge operation. Let $e \in \delta_{G}(s)$ be the hyperedge that is h-trimmed during this operation. For notational convenience, we denote $g:= e - \{s\}$ as the hyperedge formed after the h-trim operation. We define an orientation $\head_{G}:E_G\rightarrow V$ as follows: for each $e' \in E_G$,
$$\head_{G}(e') := \begin{cases}
    \head_{H}(g)& \text{ if } e' = e\\
    \head_{H}(e')& \text{ otherwise (i.e., $e'\in E_G-\{e\}$).}
\end{cases}$$  
%\knote{Replace ``otherwise'' above by ``$e'\in E_H-\{e\}$''}
\end{enumerate}
        We now show that the orientation $\overrightarrow{G} = (V, E_G, \head_G)$ is Steiner rooted $k$-hyperarc-connected. We will prove this only for the first case above, i.e. $H$ is obtained from $G$ by a h-merge almost disjoint hyperedges operation, and remark that the proof for the second case is along similar lines. Let $v \in T - r$ be a terminal. Since $\overrightarrow{H}$ is a Steiner rooted $k$-hyperarc-connected orientation of $H$, there exist $k$ hyperarc-disjoint paths $\calP := \{P_1, \ldots, P_k\}$ 
        %from $r$ to $v$ 
        from $v$ to $r$ 
        in $\overrightarrow{H}$. We will use the set of paths $\calP$ to show that there are also $k$ hyperarc-disjoint paths 
        %from $r$ to $v$ 
        from $v$ to $r$ 
        in $\overrightarrow{G}$. We recall that $e, f \in \delta_G(s)$ are the hyperedges that are h-merged and $g = e \cup f $ is the hyperedge formed during the splitting-off almost disjoint hyperedges operation. First, suppose that the hyperarc $(g, \head_H(g))$ is not present in any of the paths of $\calP$. Then, $\calP$ is a set of $k$ hyperarc-disjoint paths 
        %from $r$ to $v$ 
        from $v$ to $r$ 
        in $\overrightarrow{G}$ and the claim holds. Next, suppose that the hyperarc $(g, \head_H(g))$ is present in the path $P_j \in \calP$ for some $j \in [k]$. We note that since the paths in $\calP$ are hyperarc-disjoint, $P_j$ is the unique path containing the hyperarc $(g, \head_H(g))$ and the paths in $\calP - \{P_j\}$ are hyperarc-disjoint in $\overrightarrow{G}$. We now define a new path $\Tilde{P_j}$ 
        %from $r$ to $v$ 
        from $v$ to $r$ 
        in $\overrightarrow{G}$ that is hyperarc-disjoint from the paths in $\calP - \{P_j\}$. Let $\ell \in \Z_+$ denote the length of the path $P_j$ and let the vertices and hyperarcs along the path $P_j$ be 
        $$P_j = \left(v, \left(e_0, u_1\right), u_1, \left(e_1, u_2\right), u_2,  \ldots, u_{i-1}, \left(e_{i-1}, u_{i}\right), u_i, \left(g, u_{i+1}\right), u_{i+1}, \left(e_{i+1}, u_{i+2}\right), u_{i+1}, \ldots, u_{\ell}, \left(e_\ell, r\right), r\right).$$ 
        If the vertices $u_i, u_{i+1} \in e$, then we define the path $\Tilde{P}_j$ by replacing the hyperarc $(g, \head_H(g))=(g, u_{i+1})$ with the hyperarc $(e, \head_G(e))=(e, u_{i+1})$ in $P_j$ (the case where $u_i, u_{i+1} \in f$ can be shown similarly). 
        Suppose that $u_i \in e-f$ and $u_{i+1} \in f-e$: then, we define the path $\Tilde{P}_j$ by replacing the hyperarc $(g, \head_H(g))=(g,u_{i+1})$ with 
        %the hyperarcs $(e, s), (f, u_{i+1})$ in $P_j$ 
        the sub-path $(e, s), s, (f, u_{i+1})$ in $P_j$; we note that $\head_G(e)=s$, $s$ is a tail of $f$, and $\head_G(f)=u_{i+1}$ (the case where $u_i \in f-e$ and $u_{i+1} \in e-f$ can be shown similarly). 
        In both cases, the path $\Tilde{P}_j$ is a $v$ to $r$ path in the directed hypergraph $\overrightarrow{H}$. Moreover, the $k$ paths in $\calP - \{P_j\} + \{\Tilde{P}_j\}$ are hyperarc-disjoint. Consequently, $\overrightarrow{G}$ is a Steiner rooted $k$-hyperarc-connected orientation of $G$.
\end{proof}

\paragraph{The case of $|T|=|V|$. } \textsc{Max Steiner Rooted-Connected Orientation} problem where all vertices are terminals is polynomial-time solvable owing to a characterization for the existence of rooted $k$-hyperarc-connected orientation via weak-partition-connectivity. We define weak-partition-connectivity now. 
\begin{definition}\label{defn:weak-partition-connectivity}
    Let $k \in \Z_{\geq 0}$ be a non-negative integer. A hypergraph $H = (V, E)$ is \emph{$k$-weakly-partition-connected} if for every partition $\calP$ of $V$, we have that $$\sum_{e \in E} \left(\calP\left(e\right) - 1\right) \geq k (|\calP| - 1),$$
    where for each $e \in E$, $\calP(e)$ denotes the number of parts of $\calP$ that have non-empty intersection with the hyperedge $e$. 
    %We denote $\alpha_H := \min\{k \in \Z_{\geq 0} : H \text{ is $k$-weakly partition-connected}\}$ as the \emph{weak-partition-connectivity} of $H$. \knote{Could remove $\alpha_H$.}
\end{definition}

For a hypergraph $H=(V, E)$ and a root vertex $r\in V$, we recall that an orientation $\overrightarrow{H}=(V, E, \head: E\rightarrow V)$ of $H$ is rooted $k$-hyperarc-connected if for each $u\in V$, there exist $k$ hyperarc-disjoint paths 
%from the root vertex $r$ to the vertex $u$. 
from the vertex $u$ to the root vertex $r$. 
The next theorem is a consequence of Theorem 9.4.13 in \cite{Frank-book} and shows that weak-partition-connectivity characterizes the existence of rooted hyperarc-connected orientations. 

\begin{theorem}[Theorem 9.4.13 in \cite{Frank-book}]\label{thm:frank:weak-pc-iff-rooted-hyparc-conn-orientation}
    Let $H = (V, E)$ be an undirected hypergraph and $r \in V$ be a root vertex. There exists a rooted $k$-hyperarc-connected orientation of $H$ if and only if $H$ is $k$-weakly-partition-connected.
\end{theorem}
Theorem \ref{thm:frank:weak-pc-iff-rooted-hyparc-conn-orientation} leads to a polynomial-time algorithm to solve \textsc{Max Steiner Rooted-Connected Orientation} problem where all vertices are terminals. 
%We will now use Lemma \ref{lem:lifting-Steiner-orientation-through-splitting-off} and Theorem \ref{thm:frank:weak-pc-iff-rooted-hyparc-conn-orientation} to prove Theorems \ref{thm:rooted-steiner-orientation-in-hypergraphs} and \ref{thm:rooted-steiner-orientation-in-graphs}. 

\paragraph{Proof of Theorems \ref{thm:rooted-steiner-orientation-in-hypergraphs} and \ref{thm:rooted-steiner-orientation-in-graphs}.}
We will now use Lemma \ref{lem:lifting-Steiner-orientation-through-splitting-off} and Theorems \ref{thm:frank:weak-pc-iff-rooted-hyparc-conn-orientation} and \ref{thm:complete-splitting-off} to prove Theorems \ref{thm:rooted-steiner-orientation-in-graphs} and \ref{thm:rooted-steiner-orientation-in-hypergraphs}. 
We will prove Theorem \ref{thm:rooted-steiner-orientation-in-hypergraphs} and remark after the proof about how our proof also implies Theorem \ref{thm:rooted-steiner-orientation-in-graphs}. 
To prove Theorem \ref{thm:rooted-steiner-orientation-in-hypergraphs}, we relate connectivity and weak-partition-connectivity. We recall that a hypergraph $H=(V, E)$ is defined to be $k$-hyperedge-connected if $d_H(X)\ge k$ for every non-empty proper subset $X\subseteq V$. 
%By Menger's theorem, we know that a hypergraph $H=(V, E)$ is $k$-hyperedge-connected if and only if there exist $k$ hyperedge-disjoint paths between every pair of distinct vertices $u, v\in V$. 
It is easy to see that if a hypergraph is $k$-weakly-partition-connected, then it is also $k$-hyperedge-connected, but the converse is not necessarily true (counterexamples exist even for graphs). We show an approximate converse: if a hypergraph is $2k$-hyperedge-connected, then it is $k$-weak-partition-connected.
\begin{lemma}\label{lem:2k-conn-implies-k-weak-pc}
    %Let $H = (V, E)$ be an undirected hypergraph such that $\lambda_H$ is even. Then, $\lambda_H \leq 2 \alpha_H$.
    Let $H = (V, E)$ be an undirected hypergraph that is $2k$-hyperedge-connected for some non-negative integer $k$. Then, $H$ is $k$-weak-partition-connected.
\end{lemma}
\begin{proof}
    %The inequality $\alpha_H\leq \lambda_H$ is straightforward since the partition $\calP = (P_1, P_2)$ of $V$ where $P_1$ and $P_2$ are minimum cuts of $H$ (i.e. $d_H(P_1) = d_{H}(P_2) = \lambda_H$) satisfy $\sum_{e\in E}(\calP(e) - 1) = \lambda_H (|\calP| - 1)$. We now show the second inequality of the lemma.
%    \knote{Pick an arbitrary partition of $V$, contract the parts, and argue that the relevant ratio is at least $k$. To argue this, use the fact that the cut value of each part is at least $2k$. Such a proof would avoid the floor issue and would be simpler to follow.}
   % \shubhang{why does the proof strategy you mentioned avoid the floor issue? Perhaps having a single sequence of inequalities will make the proof easier to follow (see modification below)?}
    Let $\calP = (P_1, P_2, \ldots, P_t)$ %denote the partition of $V$ witnessing $\alpha_H$, 
    be a partition of $V$ such that $t > 1$.
    %i.e. $$\calP := \arg\min\left\{\left\lfloor \frac{\sum_{e \in E}(\calP(e) - 1)}{|\calP| - 1}\right\rfloor : \text{$\calP$ is a partition of $V$ s.t. $|\calP| > 1$}\right\}.$$
    Let $H' = (V', E')$ be the hypergraph obtained after contracting the parts of $\calP$ and discarding the singleton hyperedges. We note that $|V'| = t$ and $|e| \geq 2$ for each $e \in E'$. 
    %We first show a simple upper bound on $\lambda_H$.
    The following claim gives a useful intermediate step.
    \begin{claim} \label{claim:weak-partition-conn-vs-min-cut}
        $\frac{\sum_{e \in E'}|e|}{t} \leq \frac{2\sum_{e \in E'}(|e| - 1)}{t - 1}.$
    \end{claim}
    \begin{proof}
    First, we observe that $2t|E'| \leq \sum_{e\in E'}|e|$ because $|e| \geq 2$ for each $e \in E'$.  
    Adding $(t-1)\sum_{e \in E'}|e| - 2t|E'|$ to the LHS of the previous equation, and $(2t-1)\sum_{e \in E'}|e| - 2t|E'|$ to the RHS of the previous equation, we obtain the following inequality: $(t-1)\sum_{e \in E'}|e| \leq 2t\sum_{e \in E'}|e| - 2t|E'|$.
    This inequality implies the claim since $t>1$. 
    \end{proof}
    With the above claim, we have the following:
$$2k \leq \min_{i \in [t]}\{d_{H}(P_i)\} 
            = \min_{u \in V'}\{d_{H'}(u)\}
            \leq \frac{\sum_{u \in V'}d_{H'}(u)}{t}
            = \frac{\sum_{e \in E'}|e|}{t}
            \leq \frac{2\sum_{e \in E'}(|e| - 1)}{t - 1}
            =  \frac{2\sum_{e \in E'}(\calP(e) - 1)}{|P| - 1}.$$
    Here, the first inequality is because $H$ is $2k$-hyperedge-connected and the third inequality is by \Cref{claim:weak-partition-conn-vs-min-cut}. Thus, $H$ is $k$-weak-partition-connected.
\end{proof}

We now restate and prove \Cref{thm:rooted-steiner-orientation-in-hypergraphs}.

\thmRootedSteinerArcConnectedOrientationsHypergraphs*
\begin{proof}
    Let $u_1, \ldots, u_p$ denote an arbitrary ordering of the Steiner vertices. By repeated application of \Cref{thm:complete-splitting-off}, there exists a hypergraph $H = (V, E_H)$ that can be obtained by complete h-splitting-off vertices $u_1, \ldots, u_p$ in order while preserving local connectivity. Let $H' = (T, E_H)$ be the hypergraph obtained from $H$ by deleting the isolated vertices. We note that $H$ is Steiner $2k$-hyperedge-connected, and consequently, the hypergraph $H' = (T, E_H)$ obtained from $H$ by deleting the isolated Steiner vertices is $2k$-hyperedge-connected. By \Cref{lem:2k-conn-implies-k-weak-pc}, the hypergraph $H'$ is $k$-weak-partition-connected. By \Cref{thm:frank:weak-pc-iff-rooted-hyparc-conn-orientation}, there exists a rooted $k$-hyperarc-connected orientation $\overrightarrow{H'} = (T, E_{H}, \head_{H'}:E_H \rightarrow T)$ of the hypergraph $H'$. Consequently, the orientation $\overrightarrow{H} = (V, E_H, \head_H:E_H \rightarrow V)$, where $\head_H(e) := \head_{H'}(e)$ for each $e \in E_H$, is a Steiner rooted $k$-hyperarc-connected orientation of $H$. By \Cref{lem:lifting-Steiner-orientation-through-splitting-off} and induction on the number of h-splitting-off operations needed to obtain $H$ from $G$, we have that $G$ has a Steiner rooted $k$-hyperarc-connected orientation.
\end{proof}

\begin{remark}
Certain remarks regarding our proof of Theorem \ref{thm:rooted-steiner-orientation-in-hypergraphs} are in order.
\begin{enumerate}
\item Our proof of \Cref{thm:rooted-steiner-orientation-in-hypergraphs} is constructive and leads to a polynomial-time algorithm to find a Steiner rooted $k$-hyperarc-connected orientation of a $2k$-hyperedge-connected-hypergraph. 
\item Our proof of \Cref{thm:rooted-steiner-orientation-in-hypergraphs} also implies a proof of Theorem \ref{thm:rooted-steiner-orientation-in-graphs}: to prove Theorem \ref{thm:rooted-steiner-orientation-in-graphs}, we start from a graph $G=(V, E)$, but our local connectivity preserving complete h-splitting-off operations at Steiner vertices results in a hypergraph $H=(T, E_H)$ that is $2k$-hyperedge-connected; by Lemma \ref{lem:2k-conn-implies-k-weak-pc}, the hypergraph $H$ is $k$-weak-partition-connected;   
now Theorem \ref{thm:frank:weak-pc-iff-rooted-hyparc-conn-orientation} gives a rooted $k$-hyperarc-connected orientation of the resulting hypergraph. Such an orientation is extended to a Steiner rooted $k$-arc-connected orientation of the graph $G=(V, E)$ using Lemma \ref{lem:lifting-Steiner-orientation-through-splitting-off}. 
\item Our proof of \Cref{thm:rooted-steiner-orientation-in-hypergraphs} reveals the source of the $2$-factor gap in the approximate min-max relation of Kir\'{a}ly and Lau for \textsc{Max Steiner Rooted-Connected Orientation Problem}: it arises from the $2$-factor gap between connectivity and weak-partition-connectivity of hypergraphs. This insight immediately leads to a minimal example that exhibits tightness of the $2$-factor in Theorem %\ref{thm:rooted-steiner-orientation-in-hypergraphs} and 
\ref{thm:rooted-steiner-orientation-in-graphs}: consider the cycle graph on $3$ vertices with all vertices being terminals and an arbitrary vertex chosen as the root. This graph is $2$-edge-connected, $1$-weak-partition-connected but not $2$-weak-partition-connected, and admits a Steiner rooted $2$-arc-connected orientation but does not admit a Steiner rooted $2$-arc-connected orientation. 
%the example in Figure \ref{??} which exhibits the tightness of the $2$-factor in Theorems \ref{thm:rooted-steiner-orientation-in-hypergraphs} and \ref{thm:rooted-steiner-orientation-in-graphs}.
\end{enumerate}
\end{remark}

